\subsubsection{Analisis Mekanisme Penggunaan dan Revokasi Token}
\label{subsubsec:analisis-mekanisme-penggunaan-dan-revokasi-akses}

Pada implementasi DecMed sebelumnya, kapabilitas akses ditentukan oleh kepemilikan token dan klaim yang dibawa oleh token tersebut. Setiap operasi pembacaan maupun penulisan data RME dilakukan melalui \textit{request} yang menyertakan token akses. Komponen PRE kemudian memvalidasi token dengan mengacu pada hak akses yang tercatat pada penyimpanan \textit{on-chain} IOTA. Validasi ini digunakan untuk memastikan bahwa \textit{address} IOTA yang tercantum pada klaim token memiliki hak akses yang sesuai pada penyimpanan \textit{on-chain} dan hak akses tersebut masih berlaku. Struktur klaim token ditunjukkan pada Kode \ref{lst:decmed-jwt-claims}.

Mekanisme tersebut masih memiliki keterbatasan karena validasi berfokus pada klaim yang dibawa token, bukan pada pembuktian bahwa pengirim \textit{request} benar-benar merupakan pemilik \textit{address} yang tercantum pada klaim. Jika token akses diperoleh oleh pihak lain, misalnya melalui kebocoran atau pencurian token, pihak tersebut dapat mengirimkan \textit{request} dengan klaim \textit{address} yang masih valid. Kondisi ini membuat token dapat digunakan oleh aktor yang tidak berwenang selama klaim di dalam token masih diterima oleh sistem. Oleh karena itu, diperlukan mekanisme tambahan agar penggunaan token diikat dengan bukti identitas dari aktor yang melakukan akses.

Selain verifikasi terhadap pengguna token, mekanisme pencabutan hak akses juga perlu diperhatikan. Pada DecMed, pasien dapat mencabut hak akses sebelum masa berlaku akses berakhir. Proses pencabutan dilakukan melalui \textit{patient client} dengan mengirimkan \textit{request} revokasi ke PRE. Setelah diterima, status revokasi dicatat pada Redis dan penyimpanan \textit{on-chain} IOTA agar akses yang telah dicabut tidak lagi diterima oleh sistem.

Apabila mekanisme delegasi akses ditambahkan, sebagaimana dibutuhkan pada skenario layanan kesehatan kolaboratif di fasilitas pelayanan kesehatan pada Subbab \ref{subsubsec:analisis-skenario-layanan-klinis}, mekanisme revokasi tidak cukup hanya menangani token awal. Token hasil delegasi juga perlu diperhitungkan, termasuk hubungan antara token induk dan token turunannya. Dengan demikian, pencabutan akses dapat diterapkan tidak hanya terhadap akses yang diberikan langsung oleh pasien, tetapi juga terhadap akses yang muncul melalui proses delegasi.
